
\section{整体生成的拟凝聚层}

我们从一个结果开始，它让我们能够借助\trLineBundle{}来「够着」\trProjectiveScheme{}上的任何\trFiniteType{}的\trQCSheaf。本节通篇，$S_\bullet$ 都将是 $A$ 上一个\trFinitelyGenerated{}的\trGradedRing，在次数 $1$ 处生成。我们将证明下述结果。

\textbf{16.1.1.\ \trhTheorem{}。}\ ——\ \emph{$\Proj S_\bullet$ 上任何\trFiniteType{}的\trQCSheaf{} $\mathscr{F}$ 均能表成如下形式}
\[
\oplus_{\mathrm{finite}}\OO(-\mathfrak{m})\longrightarrow\mathscr{F}\longrightarrow 0.
\]

由于我们能够亲手摆弄这些\trLineBundle{} $\OO(-\mathfrak{m})$，这一结果将让我们对 $\Proj S_\bullet$ 上的\emph{一切}\trCoherentSheaf（特别是\trVectorBundle）都有很强的掌控。作为一个初例，它将使我们得以证明：\trProjective{} $k$-\trScheme{}上每个\trCoherentSheaf{}的\trGlobalSections{}都构成一个\trFiniteDimensional{}的\trVectorSpace（\mbox{推论~18.1.4}）。（这个事实日后将长大为这样一个事实：\trCoherentSheaf{}在\trProper{}态射下的高阶前推也是\trCoh{}的；见\mbox{定理~18.7.1}(d) 与\trGrothendieck{}凝聚定理~18.9.1。）

我们不径直奔向证明，而是拿它作个借口，来介绍一些在更宽泛的场合有用的观念（\emph{\trGlobalGeneration}（global generation）、\emph{无基点性}（base-point-freeness）、\emph{\trAmpleness}（ampleness））以及它们之间的相互关系。但首先我们拿它作借口，来提一条重要的经典结果。

\textbf{16.1.2.}\ \emph{\trHilbert{}\trSyzygy{}定理。}

对 $\mathbb{P}^n_k$ 上任一\trCoherentSheaf{} $\mathscr{F}$，\mbox{定理~16.1.1} 给出了\trSurjectiveMap{} $\phi\colon\oplus_{\mathrm{finite}}\OO(-\mathfrak{m})\to\mathscr{F}\to 0$ 的存在性。这个\trSurjectiveMap{}的\trKernel{}也是\trCoh{}的，于是迭代这一构造，我们便能构造出 $\mathscr{F}$ 的一个由\trLineBundle{}的\trDirectSum{}给出的无穷分解（infinite resolution）：
\[
\cdots\;\oplus_{\mathrm{finite}}\OO(\mathfrak{m}_{2,j})\longrightarrow\oplus_{\mathrm{finite}}\OO(\mathfrak{m}_{1,j})\longrightarrow\oplus_{\mathrm{finite}}\OO(\mathfrak{m}_{0,j})\longrightarrow\mathscr{F}\longrightarrow 0.
\]

\textbf{16.1.3.\ \trHilbert{}\trSyzygy{}定理。}\ ——\ \emph{存在一个有限的分解，其长度至多为 $n+1$。}

这一结果的最初历史，见\mbox{定理~3.6.17} 之后的评述。我们不会用到这条定理，但其证明将在选读的 \S24.4 中给出。

\textbf{16.1.4.\ \trGloballyGenerated\trSheaf。}\quad 设 $X$ 是一个\trScheme，而 $\mathscr{F}$ 是一个 $\OO$-\trModule。本节最重要的定义就是如下这个：称 $\mathscr{F}$ 为\textbf{\trGloballyGenerated}（globally generated）（或\textbf{由\trGlobalSections{}生成}（generated by global sections）），若它容许一个来自 $X$ 上某个\trFreeSheaf{}的\trSurjectiveMap：
\[
\OO^{\oplus I}\twoheadrightarrow\mathscr{F}.
\]
这里 $I$ 是某个\trIndexSet。此处所论的那些\trGlobalSections，便是 $\OO_X^{\oplus I}$ 的各个直和项里对应于 $1$ 的那 $|I|$ 个\trSection{}的像；这些像生成 $\mathscr{F}$ 的各\trStalk。若\trIndexSet{} $I$ 可以取为有限的，我们便说 $\mathscr{F}$ 是\textbf{有限整体生成的}（finitely globally generated）（或\textbf{由有限多个\trGlobalSections{}生成}（generated by a finite number of global sections））。

更多的定义，说得再细致些：我们说 $\mathscr{F}$ \textbf{在点 $p$ 处整体生成}（globally generated at a point $p$）（或有时说\textbf{在 $p$ 处由\trGlobalSections{}生成}（generated by global sections at $p$）），如果我们能找到一个 $\phi\colon\OO^{\oplus I}\to\mathscr{F}$，它在 $p$ 处的\trStalk{}上是\trSurjectiveMap，
\[
\OO_{\mathfrak{p}}^{\oplus I}\overset{\phi_{\mathfrak{p}}}{\twoheadrightarrow}\mathscr{F}_{\mathfrak{p}}.
\]
（说得更精确些，应当是：$\mathscr{F}$ 在 $p$ 处的\trStalk{}由 $\mathscr{F}$ 的\trGlobalSections{}生成。）关键的洞见在于：在 $p$ 处的整体生成，意味着 $p$ 处的每一个\trGerm{}都是诸\trGlobalSections{}的\trGerm{}的一个线性组合（组合系数取在\trLocalRing{} $\OO_{X,p}$ 中）。

注意，若 $\mathscr{F}$ 在\emph{一切}点 $p$ 处都整体生成，则 $\mathscr{F}$ 是\emph{\trGloballyGenerated}。（理由：\mbox{练习~2.4.D} 已证明\trIsomorphism{}可以在\trStalk{}的层面上检验；同一论证的一个更容易的版本表明，\trSurjectiveMap{}性也可以在\trStalk{}的层面上检验。）请注意，我们可以为整个 $X$ 取定单一的\trIndexSet，办法是把各个 $p$ 的诸\trIndexSet{}取并。

\textbf{16.1.A.（容易的练习）}\quad $X$ 上的一个\trInvertibleSheaf{} $\mathscr{L}$ 是\trGloballyGenerated，当且仅当对任意点 $p\in X$，都存在 $\mathscr{L}$ 的一个在 $p$ 处不消失的\trSection；也就是说，$\mathscr{L}$ 是\trBasePointFree{}的（\mbox{定义~15.2.1}）。

\textbf{16.1.B.（容易的练习：现实检验）}\quad 证明：每个\trAffineScheme{}上的每个\trQCSheaf{}都是\trGloballyGenerated。证明：每个\trAffineScheme{}上的每个\trFiniteType{}的\trQCSheaf{}都由有限多个\trGlobalSections{}生成。（两者共用的提示：对任何 $A$-\trModule{} $M$，总有一个从自由 $A$-\trModule{}到 $M$ 的\trSurjectiveMap。）

显然，若 $\mathscr{F}$ 与 $\mathscr{G}$ 都\trGloballyGenerated，则 $\mathscr{F}\oplus\mathscr{G}$ 亦然（对有限整体生成、在点 $p$ 处生成等等情形亦复如是）。对张量积也有类似的道理：

\textbf{16.1.C.（容易的练习：整体生成的$\otimes$整体生成的是整体生成的）}\quad 证明：若\trQCSheaf{} $\mathscr{F}$ 与 $\mathscr{G}$ 在点 $p$ 处整体生成，则 $\mathscr{F}\otimes\mathscr{G}$ 亦然。（若 $\mathscr{F}$ 与 $\mathscr{G}$ 是\trInvertibleSheaf，这正是\mbox{练习~15.2.B}。）

\textbf{16.1.D.（容易但重要的练习）}\quad 设 $\mathscr{F}$ 是 $X$ 上一个\trFiniteType{}的\trQCSheaf。

（a）证明：$\mathscr{F}$ 在点 $p$ 处整体生成，当且仅当「$\mathscr{F}$ 的\trFiber{}在 $p$ 处由\trGlobalSections{}生成」，也就是说，诸\trGlobalSections{}作为一个\trVectorSpace{}张成\trFiber{} $\mathscr{F}_{\mathfrak{p}}/\mathfrak{m}\mathscr{F}_{\mathfrak{p}}$，其中 $\mathfrak{m}$ 是 $\OO_{X,\mathfrak{p}}$ 的\trMaximalIdeal。（提示：几何\trNakayama{}引理，\mbox{练习~14.3.D}。）

（b）证明：若 $\mathscr{F}$ 在点 $p$ 处整体生成，则「$\mathscr{F}$ 在 $p$ 附近整体生成」：存在 $p$ 的一个\trOpen\trNeighborhood{} $\mathrm{U}$，使得 $\mathscr{F}$ 在 $\mathrm{U}$ 的每一点处都整体生成。

（c）进而设 $X$ 是一个\trQuasicompact{}的\trScheme。证明：若 $\mathscr{F}$ 在 $X$ 的一切\trClosedPoint{}处都整体生成，则 $\mathscr{F}$ 在 $X$ 的一切点处都整体生成。（注意，非空的\trQuasicompact{}概形\emph{确有}\trClosedPoint，见\mbox{练习~5.1.E}。）

\textbf{16.1.E.（容易的练习）}\quad 若 $\mathscr{F}$ 是 $X$ 上一个\trFiniteType{}的\trQCSheaf，而 $X$ 是\trQuasicompact{}的，证明：$\mathscr{F}$ 是\trGloballyGenerated，当且仅当它由\emph{有限多个}\trGlobalSections{}生成。

现在我们能够陈述\trSerre{}的一条著名结果了。

\textbf{16.1.5.\ \trSerre{}定理~A。}\ ——\ \emph{设 $S_\bullet$ 在次数 $1$ 处生成，并且在 $A=S_0$ 上\trFinitelyGenerated。令 $\mathscr{F}$ 为 $\Proj S_\bullet$ 上任一\trFiniteType{}的\trQCSheaf。那么存在某个 $\mathfrak{m}_0$，使得对一切 $\mathfrak{m}\ge\mathfrak{m}_0$，$\mathscr{F}(\mathfrak{m})$ 都是有限整体生成的。}

我们本可以此刻就径直证明\trSerre{}定理~A，但我们仍将继续拿它当借口来引入更多的想法；它将是\mbox{定理~16.2.2} 的一个推论。

\textbf{16.1.6.}\ \emph{\mbox{定理~16.1.1} 在\trSerre{}定理~A（\mbox{定理~16.1.5}）假设下的证明。}\quad 设我们有 $\mathscr{F}(\mathfrak{m})$ 的 $n$ 个\trGlobalSections{} $s_1,\dots,s_n$，它们生成 $\mathscr{F}(\mathfrak{m})$。这就给出一个映射
\[
\OO^{\oplus n}_{\Proj S_\bullet}\longrightarrow\mathscr{F}(\mathfrak{m})
\]
它在任何\trOpenSet{}上都由 $(f_1,\dots,f_n)\mapsto f_1s_1+\cdots+f_ns_n$ 给出。由于这些\trGlobalSections{}生成 $\mathscr{F}(\mathfrak{m})$，这是一个\trSurjectiveMap。用 $\OO(-\mathfrak{m})$ 作张量（它是正合的，因为与任何\trLocallyFreeSheaf{}作张量都正合，\mbox{练习~14.1.F}）便给出所要的结果。\hfill$\square$
